\documentclass[xcolor=dvipsnames]{beamer}
\usepackage[utf8]{inputenc}
\usepackage[german]{babel}
\usepackage{xcolor}
\usepackage{graphicx}
\usepackage{MnSymbol}
\usepackage{stmaryrd}
\usepackage{colortbl}
\usepackage{caption}
\usepackage{comment}
\usepackage{pdfpages}
\usepackage{listings}
\usepackage{color}
\usepackage{booktabs}
\usepackage{soul}
\usepackage[normalem]{ulem}
\usepackage{tcolorbox}
\usepackage{amsmath}
\usepackage{amssymb}

%%%%%%%%%%%%%%%%%%%%%%%%%%%%%%%%%%%%%%%%%%%%%%%%%
\usepackage{pgf}
\usepackage{etex}
\usepackage{tikz,pgfplots}
\usetikzlibrary{calc,trees,positioning,arrows,chains,shapes.geometric,decorations.pathreplacing,decorations.pathmorphing,shapes,matrix,shapes.symbols}

\usetheme{Madrid}
\usecolortheme{dolphin}
\usefonttheme{professionalfonts}
\useoutertheme{infolines}
\useinnertheme{circles}

\newtheorem*{bem}{Bemerkung}

%%%%%%%%%%%%%%%%%%%%%%%%%%%%%%%%%%%%%%%%%%%%%%%%%
% Code Listing Settings
%%%%%%%%%%%%%%%%%%%%%%%%%%%%%%%%%%%%%%%%%%%%%%%%%
\definecolor{dkgreen}{rgb}{0,0.6,0}
\definecolor{gray}{rgb}{0.5,0.5,0.5}
\definecolor{mauve}{rgb}{0.58,0,0.82}
\definecolor{codebg}{rgb}{0.95,0.95,0.95}

\lstset{
  language=C++,
  aboveskip=3mm,
  belowskip=3mm,
  showstringspaces=false,
  columns=flexible,
  basicstyle={\small\ttfamily},
  numbers=left,
  numberstyle=\tiny\color{gray},
  keywordstyle=\color{blue},
  commentstyle=\color{dkgreen},
  stringstyle=\color{mauve},
  breaklines=true,
  breakatwhitespace=true,
  tabsize=2,
  backgroundcolor=\color{codebg},
  frame=single
}

%%%%%%%%%%%%%%%%%%%%%%%%%%%%%%%%%%%%%%%%%%%%%%%%%

\title[N-Body Simulation]{N-Body Simulation mit Barnes-Hut Algorithmus}
\author{Ihr Name}
\institute{TU Berlin}
\date{\today}

\begin{document}

% Title Page
\begin{frame}
  \titlepage
\end{frame}

% Table of Contents
\begin{frame}
\frametitle{Inhaltsverzeichnis}
\tableofcontents
\end{frame}

\begin{frame}{Benötigte Screenshots - Übersicht}
\begin{alertblock}{Visualisierungen für die Präsentation}
Folgende Screenshots werden benötigt:
\end{alertblock}

\begin{enumerate}
    \item \textbf{QuadTree Aufbau (2D)}
    \begin{itemize}
        \item Schrittweiser Aufbau, rekursive Unterteilung
    \end{itemize}

    \item \textbf{Octree Struktur (3D)}
    \begin{itemize}
        \item 3D-Würfel mit hierarchischer Struktur
    \end{itemize}

    \item \textbf{Partikel Einfügen}
    \begin{itemize}
        \item Animation des rekursiven Einfügens
    \end{itemize}

    \item \textbf{Massenschwerpunkt-Propagation}
    \begin{itemize}
        \item Bottom-Up Berechnung mit Farbcodierung
    \end{itemize}

    \item \textbf{Kraft-Berechnung mit Theta}
    \begin{itemize}
        \item Traversierung visualisiert, grün/rot markiert
    \end{itemize}
\end{enumerate}
\end{frame}

%%%%%%%%%%%%%%%%%%%%%%%%%%%%%%%%%%%%%%%%%%%%%%%%%
\section{Einleitung}
%%%%%%%%%%%%%%%%%%%%%%%%%%%%%%%%%%%%%%%%%%%%%%%%%

\begin{frame}{Das N-Körper-Problem}
\begin{block}{Problemstellung}
Simulation von N Partikeln, die sich gegenseitig gravitativ anziehen
\end{block}

\begin{columns}
\column{0.5\textwidth}
\textbf{Input:}
\begin{itemize}
    \item N Partikel
    \item Position $(x, y, z)$
    \item Geschwindigkeit $(v_x, v_y, v_z)$
    \item Masse $m$
\end{itemize}

\column{0.5\textwidth}
\textbf{Output:}
\begin{itemize}
    \item Neue Positionen
    \item Neue Geschwindigkeiten
    \item Nach Zeitschritt $\Delta t$
\end{itemize}
\end{columns}

\vspace{0.5cm}
\begin{exampleblock}{Anwendungen}
Galaxiensimulation, Sternentstehung, Planetenbewegung, Molekulardynamik
\end{exampleblock}
\end{frame}

%%%%%%%%%%%%%%%%%%%%%%%%%%%%%%%%%%%%%%%%%%%%%%%%%
\section{Naive Lösung: Brute Force}
%%%%%%%%%%%%%%%%%%%%%%%%%%%%%%%%%%%%%%%%%%%%%%%%%

\begin{frame}[fragile]{Brute Force Algorithmus}
\begin{block}{Direkter Ansatz}
Berechne Kraft zwischen \textbf{jedem Partikelpaar}
\end{block}

\begin{lstlisting}
for (int i = 0; i < N; i++) {
    for (int j = 0; j < N; j++) {
        if (i != j) {
            Vector3 r = particles[j].pos - particles[i].pos;
            double dist = length(r);
            double force = G * particles[i].mass * particles[j].mass
                         / (dist*dist + eps);
            particles[i].acceleration += force * normalize(r);
        }
    }
}
\end{lstlisting}
\end{frame}

\begin{frame}{Brute Force: Komplexität}
\begin{block}{Zeitkomplexität: $O(n^2)$}
Für N Partikel: $\frac{N \cdot (N-1)}{2}$ Berechnungen
\end{block}

\begin{table}
\centering
\begin{tabular}{r|r}
\textbf{Partikel} & \textbf{Berechnungen/Frame} \\ \hline
100 & 4.950 \\
1.000 & 499.500 \\
10.000 & 49.995.000 \\
100.000 & $\sim$5 Milliarden
\end{tabular}
\end{table}

\begin{alertblock}{Problem}
Wird \textbf{extrem schnell} sehr langsam! \\
Für große N praktisch unmöglich.
\end{alertblock}
\end{frame}

%%%%%%%%%%%%%%%%%%%%%%%%%%%%%%%%%%%%%%%%%%%%%%%%%
\section{Barnes-Hut Algorithmus}
%%%%%%%%%%%%%%%%%%%%%%%%%%%%%%%%%%%%%%%%%%%%%%%%%

\begin{frame}{Barnes-Hut: Kernidee}
\begin{block}{Problem mit Brute Force}
$O(n^2)$ - zu langsam für große N!
\end{block}

\begin{block}{Lösung: Approximation}
Ferne Partikelgruppen als \textbf{einzelner Massenpunkt} behandeln!
\end{block}

\begin{center}
\begin{tikzpicture}[scale=0.8]
% Particle being calculated
\fill[red] (0,0) circle (3pt);
\node[below] at (0,-0.3) {\textcolor{red}{Partikel i}};

% Far away group
\draw[thick,blue] (6,0) circle (1.5cm);
\fill[blue] (6,0.3) circle (1.5pt);
\fill[blue] (6.5,0) circle (1.5pt);
\fill[blue] (5.8,-0.4) circle (1.5pt);
\fill[blue] (6.2,0.5) circle (1.5pt);
\fill[blue] (5.5,0.2) circle (1.5pt);

% Arrow and approximation
\draw[->,very thick] (1,0.5) -- (4.5,0.5) node[midway,above] {weit genug?};
\draw[->,very thick,green!50!black] (6,0) -- (7.8,0);
\fill[green!50!black] (8.2,0) circle (4pt);
\node[below] at (8.2,-0.3) {\textcolor{green!50!black}{$\approx$ 1 Punkt}};
\end{tikzpicture}
\end{center}

\begin{exampleblock}{Resultat}
$O(n^2) \rightarrow O(n \log n)$ --- 30-250x Speedup!
\end{exampleblock}
\end{frame}

\begin{frame}{Barnes-Hut: 3 Schritte}
\begin{center}
\large
\textbf{Wie funktioniert Barnes-Hut?}
\end{center}

\begin{enumerate}
    \item \textbf{Schritt 1}: Octree aufbauen durch Partikel-Einfügen
    \begin{itemize}
        \item Baum entsteht rekursiv beim Einfügen jedes Partikels
        \item Subdivision bei $>1$ Partikel pro Knoten
    \end{itemize}
    \item \textbf{Schritt 2}: Massenschwerpunkte berechnen
    \item \textbf{Schritt 3}: Kräfte mit Theta-Kriterium berechnen
\end{enumerate}

\vspace{0.5cm}
\begin{alertblock}{Wichtig}
Wir erklären jeden Schritt \textbf{im Detail} mit Visualisierungen!
\end{alertblock}
\end{frame}

%%%%%%%%%%%%%%%%%%%%%%%%%%%%%%%%%%%%%%%%%%%%%%%%%
\section{Schritt 1: Octree aufbauen durch Partikel-Einfügen}
%%%%%%%%%%%%%%%%%%%%%%%%%%%%%%%%%%%%%%%%%%%%%%%%%

\begin{frame}{Schritt 1.1: QuadTree (2D) - Konzept}
\begin{columns}
\column{0.5\textwidth}
\textbf{Hierarchische Raumpartitionierung}
\begin{itemize}
    \item 2D-Raum rekursiv unterteilen
    \item Jeder Knoten hat 4 Kinder
    \item Stopp bei $\leq 1$ Partikel
\end{itemize}

\vspace{0.3cm}
\textbf{4 Quadranten:}
\begin{itemize}
    \item Nordwest (NW)
    \item Nordost (NO)
    \item Südwest (SW)
    \item Südost (SO)
\end{itemize}

\column{0.5\textwidth}
\begin{center}
\begin{tikzpicture}[scale=0.7]
% Main square
\draw[very thick] (0,0) rectangle (4,4);

% Level 1 divisions
\draw[thick] (2,0) -- (2,4);
\draw[thick] (0,2) -- (4,2);

% Level 2 divisions (SW quadrant)
\draw (0,1) -- (2,1);
\draw (1,0) -- (1,2);

% Particles
\fill[red] (0.5,0.5) circle (2pt);
\fill[red] (1.5,0.5) circle (2pt);
\fill[red] (0.5,1.5) circle (2pt);
\fill[red] (3,3) circle (2pt);
\fill[red] (2.8,1.2) circle (2pt);

% Labels
\node at (1,3) {NW};
\node at (3,3) {NO};
\node at (1,1) {SW};
\node at (3,1) {SO};
\end{tikzpicture}
\end{center}
\end{columns}

\begin{exampleblock}{Warum 2D zuerst?}
2D QuadTree ist einfacher zu visualisieren - gleiches Prinzip wie 3D Octree!
\end{exampleblock}
\end{frame}

\begin{frame}{QuadTree: Partikel einfügen}
\begin{block}{Algorithmus}
\textbf{So entsteht der Baum:}
\end{block}

\begin{enumerate}
    \item \textbf{Finde Bounding Box}: Min/Max x, y aller Partikel
    \item \textbf{Erstelle Wurzel}: Knoten der alle Partikel enthält
    \item \textbf{Rekursives Einfügen}:
    \begin{itemize}
        \item Falls Knoten leer $\rightarrow$ Partikel direkt einfügen
        \item Falls $\leq 1$ Partikel $\rightarrow$ Blattknoten, fertig!
        \item Falls $> 1$ Partikel $\rightarrow$ Unterteile in \textbf{4 Quadranten}
        \item Bestimme Quadrant des Partikels (x, y Vergleich)
        \item Rekursiere in korrekten Quadranten
    \end{itemize}
\end{enumerate}
\end{frame}

\begin{frame}[fragile]{QuadTree-Aufbau: Pseudocode}
\begin{lstlisting}[language=Python, basicstyle=\footnotesize\ttfamily]
function subdivide(node):
    // Basisfall: <= 1 Partikel -> Fertig
    if node.particleCount <= 1:
        node.isLeaf = true
        return

    // Erstelle 4 Kinder (halbe Grosse)
    for quadrant = 0 to 3:
        child = createChild()
        child.size = node.size / 2
        child.center = berechnePosition(quadrant, node.center)

    // Verteile Partikel auf Kinder
    for particle in node:
        quadrant = findQuadrant(particle, node.center)
        children[quadrant].add(particle)

    // Rekursion: Wiederhole fur alle Kinder
    for child in children:
        subdivide(child)
\end{lstlisting}

\textbf{Kernidee:} Teile Raum in 4 Teile, verteile Partikel, rekursiere!
\end{frame}

\begin{frame}[fragile]{Quadranten-Berechnung}
\begin{block}{Zu welchem Quadranten gehört ein Partikel?}
Vergleiche Position mit Knoten-Zentrum
\end{block}

\begin{lstlisting}[basicstyle=\small\ttfamily]
int getQuadrant(Particle p, Vector2 center) {
    int quadrant = 0;
    if (p.x >= center.x) quadrant |= 1;  // x-Richtung
    if (p.y >= center.y) quadrant |= 2;  // y-Richtung
    return quadrant;  // 0-3
}
\end{lstlisting}

\begin{center}
\begin{tabular}{c|cc|l}
\textbf{Quadrant} & \textbf{x} & \textbf{y} & \textbf{Bezeichnung} \\ \hline
0 & - & - & SW (Südwest) \\
1 & + & - & SO (Südost) \\
2 & - & + & NW (Nordwest) \\
3 & + & + & NO (Nordost)
\end{tabular}
\end{center}
\end{frame}

\begin{frame}[fragile]{QuadTree: Datenstruktur}
\begin{lstlisting}[basicstyle=\small\ttfamily]
struct QuadTreeNode {
    // Bounding Box
    Vector2 center;
    double size;

    // Massendaten
    double mass;
    Vector2 centerOfMass;  // COM

    // Baumstruktur
    int firstChild;     // Index des ersten Kinds
                        // Kinder: [firstChild + quadrant]
                        // quadrant = 0-3 (berechnet on-the-fly)
    int particleIndex;  // Falls Blatt: Partikel-ID
    bool isLeaf;
};
\end{lstlisting}

\begin{itemize}
    \item \textbf{Interne Knoten}: 4 Kinder sequenziell gespeichert
    \item \textbf{Blattknoten}: Enthalten $\leq 1$ Partikel
    \item \textbf{Zugriff}: \texttt{children[firstChild + quadrant]}
\end{itemize}
\end{frame}

\begin{frame}{QuadTree Aufbau - Demo}
\begin{columns}
\column{0.6\textwidth}
\begin{center}
\only<1>{\includegraphics[width=\textwidth]{assets/quadtree_build1.png}}
\only<2>{\includegraphics[width=\textwidth]{assets/quadtree_build2.png}}
\only<3>{\includegraphics[width=\textwidth]{assets/quadtree_build3.png}}
\only<4>{\includegraphics[width=\textwidth]{assets/quadtree_build4.png}}
\only<5>{\includegraphics[width=\textwidth]{assets/quadtree_build5.png}}
\only<6>{\includegraphics[width=\textwidth]{assets/quadtree_build6.png}}
\end{center}

\column{0.4\textwidth}
\textbf{Wichtige Regeln:}
\begin{itemize}
    \item Jedes Quad kann nur 1 Körper enthalten
    \item Bei $>1$ Körper: Quad wird geteilt
    \item Start mit Bounding Box aller Partikel
    \item Rekursive Unterteilung in 4 Quadranten
    \item Stopp bei $\leq 1$ Partikel pro Quad
    \item Erzeugt hierarchische Baumstruktur
\end{itemize}
\end{columns}
\end{frame}

\begin{frame}[fragile]{Octree: Von 2D zu 3D}
\begin{columns}
\column{0.5\textwidth}
\textbf{Der einzige Unterschied:}

\vspace{0.3cm}
\begin{center}
\begin{tabular}{l|c|c}
& \textbf{QuadTree} & \textbf{Octree} \\ \hline
Dimensionen & 2D (x, y) & 3D (x, y, z) \\
Kinder & \textbf{4} & \textbf{8} \\
Unterteilung & 4 Quadranten & 8 Oktanten \\
\end{tabular}
\end{center}

\vspace{0.3cm}
\textbf{Oktanten-Berechnung (binär):}
\begin{lstlisting}[basicstyle=\tiny\ttfamily]
int octant = 0;
if (p.x >= center.x) octant |= 1;
if (p.y >= center.y) octant |= 2;
if (p.z >= center.z) octant |= 4;
// octant = 0-7
\end{lstlisting}
\small Bitweise OR-Operation: \\
\textcolor{blue}{z-Achse (4) + y-Achse (2) + x-Achse (1)}

\column{0.5\textwidth}
\begin{center}
\begin{tikzpicture}[scale=0.6]
% Draw cube
\coordinate (O) at (0,0,0);
\coordinate (A) at (3,0,0);
\coordinate (B) at (3,3,0);
\coordinate (C) at (0,3,0);
\coordinate (D) at (0,0,3);
\coordinate (E) at (3,0,3);
\coordinate (F) at (3,3,3);
\coordinate (G) at (0,3,3);

% Draw edges
\draw[thick] (O) -- (A) -- (B) -- (C) -- cycle;
\draw[thick] (D) -- (E) -- (F) -- (G) -- cycle;
\draw[thick] (O) -- (D);
\draw[thick] (A) -- (E);
\draw[thick] (B) -- (F);
\draw[thick] (C) -- (G);

% Draw internal divisions
\draw[thick,blue] (1.5,0,0) -- (1.5,3,0) -- (1.5,3,3) -- (1.5,0,3) -- cycle;
\draw[thick,blue] (0,1.5,0) -- (3,1.5,0) -- (3,1.5,3) -- (0,1.5,3) -- cycle;
\draw[thick,blue] (0,0,1.5) -- (3,0,1.5) -- (3,3,1.5) -- (0,3,1.5) -- cycle;

% Particles
\fill[red] (0.5,0.5,0.5) circle (2pt);
\fill[red] (2.5,2.5,2.5) circle (2pt);
\end{tikzpicture}

\vspace{0.3cm}
\small $2 \times 2 \times 2 = 8$ Oktanten
\end{center}
\end{columns}
\end{frame}

\begin{frame}{Octree - 3D Demo}
\begin{center}
\vspace{2cm}
{\Huge \textbf{LIVE DEMO}}

\vspace{1cm}
{\Large 3D Octree Visualisierung}

\vspace{0.5cm}
\begin{itemize}
    \item Interaktive 3D-Darstellung
    \item Hierarchische Baumstruktur
    \item Schrittweises Einfügen von Partikeln
\end{itemize}
\end{center}
\end{frame}

%%%%%%%%%%%%%%%%%%%%%%%%%%%%%%%%%%%%%%%%%%%%%%%%%
\section{Schritt 2: Massenschwerpunkt berechnen}
%%%%%%%%%%%%%%%%%%%%%%%%%%%%%%%%%%%%%%%%%%%%%%%%%

\begin{frame}{Schritt 2: Massenschwerpunkt (Center of Mass)}
\begin{block}{Warum wichtig?}
Erlaubt uns, ganze Regionen als \textbf{einzelnen Punkt} zu behandeln
\end{block}

\begin{alertblock}{Wann wird berechnet?}
COM wird \textbf{während des Baum-Aufbaus} berechnet! \\
Nach jeder Partikel-Einfügung wird COM inkrementell aktualisiert.
\end{alertblock}

\begin{block}{Bottom-Up Berechnung}
Von Blättern zur Wurzel:
\end{block}

\textbf{Blattknoten:}
\begin{equation}
\text{COM} = \text{Partikelposition}, \quad m = \text{Partikelmasse}
\end{equation}

\textbf{Interne Knoten:}
\begin{align}
m_{\text{total}} &= \sum_{i=1}^{8} m_{\text{child},i} \\
\text{COM}_x &= \frac{\sum_{i=1}^{8} m_{\text{child},i} \cdot \text{COM}_{x,i}}{m_{\text{total}}}
\end{align}
(Analog für y, z)
\end{frame}

\begin{frame}[fragile]{Massenschwerpunkt: Pseudocode}
\begin{lstlisting}[language=Python, basicstyle=\footnotesize\ttfamily]
function updateCenterOfMass(node):
    if node.isLeaf:
        node.mass = particle.mass
        node.COM = particle.position
    else:
        // Rekursion: Kinder zuerst (bottom-up)
        for i = 0 to 3:
            updateCenterOfMass(children[node.firstChild + i])

        // Summiere 4 Kinder
        totalMass = 0
        weightedPos = (0, 0)
        for i = 0 to 3:
            totalMass += child[i].mass
            weightedPos += child[i].mass * child[i].COM

        node.mass = totalMass
        node.COM = weightedPos / totalMass
\end{lstlisting}
\end{frame}

\begin{frame}{Schritt 2: Visualisierung - Massenschwerpunkt-Propagation}
\begin{center}
\textit{Baum-Aufbau \hspace{3cm} COM Propagation}

\vspace{0.3cm}
\begin{columns}
\column{0.5\textwidth}
\begin{center}
\only<1>{\includegraphics[width=\textwidth]{assets/quadtree_build1.png}}
\only<2>{\includegraphics[width=\textwidth]{assets/quadtree_build2.png}}
\only<3>{\includegraphics[width=\textwidth]{assets/quadtree_build3.png}}
\only<4>{\includegraphics[width=\textwidth]{assets/quadtree_build4.png}}
\only<5>{\includegraphics[width=\textwidth]{assets/quadtree_build5.png}}
\only<6>{\includegraphics[width=\textwidth]{assets/quadtree_build6.png}}
\end{center}

\column{0.5\textwidth}
\begin{center}
\only<1>{\includegraphics[width=\textwidth]{assets/propagation1.png}}
\only<2>{\includegraphics[width=\textwidth]{assets/propagation2.png}}
\only<3>{\includegraphics[width=\textwidth]{assets/propagation3.png}}
\only<4>{\includegraphics[width=\textwidth]{assets/propagation4.png}}
\only<5>{\includegraphics[width=\textwidth]{assets/propagation5.png}}
\only<6>{\includegraphics[width=\textwidth]{assets/propagation6.png}}
\end{center}
\end{columns}

\vspace{0.3cm}
\only<1>{\small Selbe Baum-Struktur: Links zeigt räumliche Aufteilung, rechts zeigt Massen-Propagation}
\end{center}
\end{frame}

%%%%%%%%%%%%%%%%%%%%%%%%%%%%%%%%%%%%%%%%%%%%%%%%%
\section{Schritt 3: Kräfte berechnen}
%%%%%%%%%%%%%%%%%%%%%%%%%%%%%%%%%%%%%%%%%%%%%%%%%

\begin{frame}{Theta-Kriterium ($\theta$-criterion)}
\begin{block}{Entscheidungsregel}
\begin{equation}
\frac{s}{d} < \theta \quad \Rightarrow \quad \text{Approximiere als Punktmasse}
\end{equation}
\end{block}

\begin{itemize}
    \item $s$: Größe des Octree-Knotens (Kantenlänge)
    \item $d$: Abstand zwischen Partikel und Knoten-COM
    \item $\theta$: Genauigkeitsparameter (typisch 0.5 - 1.0)
\end{itemize}

\vspace{0.3cm}
\begin{center}
\begin{tikzpicture}[scale=0.8]
% Particle
\fill[red] (0,0) circle (3pt);
\node[below] at (0,-0.3) {Partikel};

% Node
\draw[thick,blue] (5,0) rectangle (7,2);
\fill[blue] (6,1) circle (2pt);
\node[above] at (6,2.2) {Knoten (COM)};

% Distance
\draw[<->,very thick] (0.2,-0.8) -- (5.8,-0.8) node[midway,below] {$d$};

% Size
\draw[<->,very thick] (7.3,0) -- (7.3,2) node[midway,right] {$s$};
\end{tikzpicture}
\end{center}
\end{frame}

\begin{frame}{Theta-Parameter: Trade-off}
\begin{table}
\centering
\begin{tabular}{c|l|l}
\textbf{$\theta$} & \textbf{Genauigkeit} & \textbf{Performance} \\ \hline
0.0 & Exakt (wie Brute Force) & Sehr langsam \\
0.3 & Sehr hoch & Langsam \\
\textcolor{blue}{\textbf{0.7}} & \textcolor{blue}{\textbf{Gut}} & \textcolor{blue}{\textbf{Gut (Standard)}} \\
1.0 & Mittel & Schnell \\
2.0 & Niedrig & Sehr schnell
\end{tabular}
\end{table}

\begin{alertblock}{Geometrische Bedeutung}
$\theta$ kontrolliert den Winkel, unter dem ein Knoten erscheint:
\begin{itemize}
    \item Kleines $\theta$ $\rightarrow$ Knoten muss sehr klein erscheinen
    \item Großes $\theta$ $\rightarrow$ Mehr Approximation erlaubt
\end{itemize}
\end{alertblock}
\end{frame}

\begin{frame}[fragile]{Kraft-Berechnung: Code}
\begin{lstlisting}[basicstyle=\tiny\ttfamily]
void calculateForce(Particle& p, OctreeNode& node, double theta) {
    if (node.mass == 0) return;

    // Abstand zum Massenschwerpunkt
    Vector3 r = node.COM - p.position;
    double dist = length(r);

    // Theta-Kriterium
    if (node.isLeaf || (node.size / dist < theta)) {
        // APPROXIMIERE als Punktmasse
        double force = G * p.mass * node.mass /
                       (dist*dist + softening);
        p.acceleration += force * normalize(r);
    } else {
        // Zu nah -> Rekursiere in Kinder
        for (int i = 0; i < 8; i++) {
            calculateForce(p, children[node.firstChild + i], theta);
        }
    }
}
\end{lstlisting}
\end{frame}

\begin{frame}{Schritt 3: Barnes-Hut Traversierung}
\begin{center}
\vspace{1.5cm}
{\Huge \textbf{LIVE DEMO}}

\vspace{1cm}
{\Large Barnes-Hut Kraft-Berechnung}

\vspace{0.5cm}
\begin{itemize}
    \item Traversierung visualisiert
    \item Start bei Wurzel
    \item Theta-Kriterium in Aktion
    \item \textcolor{green!50!black}{Grün}: Approximiert (far enough)
    \item \textcolor{red}{Rot}: Zu nah (recurse deeper)
    \item Vergleich: Verschiedene $\theta$-Werte
\end{itemize}
\end{center}
\end{frame}

%%%%%%%%%%%%%%%%%%%%%%%%%%%%%%%%%%%%%%%%%%%%%%%%%
\section{Zeitintegration}
%%%%%%%%%%%%%%%%%%%%%%%%%%%%%%%%%%%%%%%%%%%%%%%%%

\begin{frame}{Zeitintegration - Position \& Geschwindigkeit aktualisieren}
\begin{block}{Aktualisiere Positionen und Geschwindigkeiten}
Nach Kraft-Berechnung: Bewege Partikel vorwärts in der Zeit
\end{block}

\textbf{Euler-Methode} (einfach):
\begin{align}
\vec{v}_{\text{neu}} &= \vec{v}_{\text{alt}} + \vec{a} \cdot \Delta t \\
\vec{x}_{\text{neu}} &= \vec{x}_{\text{alt}} + \vec{v}_{\text{neu}} \cdot \Delta t
\end{align}

\textbf{Leapfrog/Verlet} (genauer, energieerhaltend):
\begin{align}
\vec{v}_{1/2} &= \vec{v} + \vec{a} \cdot \frac{\Delta t}{2} \\
\vec{x}_{\text{neu}} &= \vec{x} + \vec{v}_{1/2} \cdot \Delta t \\
\vec{a}_{\text{neu}} &= \text{calculateAcceleration}(\vec{x}_{\text{neu}}) \\
\vec{v}_{\text{neu}} &= \vec{v}_{1/2} + \vec{a}_{\text{neu}} \cdot \frac{\Delta t}{2}
\end{align}
\end{frame}

%%%%%%%%%%%%%%%%%%%%%%%%%%%%%%%%%%%%%%%%%%%%%%%%%
\section{Edge Cases \& Optimierungen}
%%%%%%%%%%%%%%%%%%%%%%%%%%%%%%%%%%%%%%%%%%%%%%%%%

\begin{frame}{Edge Cases - Spezialfälle behandeln}
\begin{block}{Was passiert bei...}
\end{block}

\begin{itemize}
    \item \textbf{Partikeln an derselben Position?}
    \begin{itemize}
        \item Softening Parameter $\varepsilon$ verhindert Division durch Null
        \item $r \rightarrow \sqrt{r^2 + \varepsilon^2}$
    \end{itemize}

    \vspace{0.3cm}
    \item \textbf{Sehr ungleichen Massenverteilungen?}
    \begin{itemize}
        \item Theta-Parameter anpassen
        \item Kleineres $\theta$ für höhere Genauigkeit
    \end{itemize}

    \vspace{0.3cm}
    \item \textbf{Leeren Oktanten?}
    \begin{itemize}
        \item Einfach überspringen (mass = 0 check)
        \item Kein unnötiger Speicher
    \end{itemize}
\end{itemize}
\end{frame}

\begin{frame}[fragile]{Implementierungs-Optimierungen}
\begin{block}{1. Flat Vector Storage}
\begin{itemize}
    \item Octree-Knoten in \texttt{std::vector} speichern
    \item Indizes statt Pointer verwenden
    \item \textbf{Vorteil}: Bessere Cache-Lokalität (CPU-Cache!)
\end{itemize}
\end{block}

\begin{block}{2. Implizite Child-Indexierung}
\begin{lstlisting}
// Kinder werden sequenziell gespeichert
node.firstChild = nextIndex;
// Kind i bei: nodes[firstChild + i]  (i = 0..7)
\end{lstlisting}
\textbf{Vorteil}: Kein Array von 8 Pointern, weniger Memory-Zugriffe
\end{block}

\begin{alertblock}{Warum ist das wichtig?}
Cache-Miss ist $\sim$100x langsamer als Cache-Hit! \\
Sequenzieller Speicher $\rightarrow$ bessere Performance
\end{alertblock}
\end{frame}

\begin{frame}{Tree Rebuild Interval}
\begin{block}{Octree nicht jeden Frame neu aufbauen}
\begin{itemize}
    \item Partikel bewegen sich nur wenig pro Zeitschritt
    \item Octree-Aufbau ist teuer
    \item \textbf{Lösung}: Rebuild nur alle N Schritte (z.B. N=5)
\end{itemize}
\end{block}

\begin{table}
\centering
\begin{tabular}{c|l}
\textbf{Rebuild Interval} & \textbf{Effekt} \\ \hline
1 & Höchste Genauigkeit, langsam \\
5 & Guter Trade-off \\
10 & Schneller, weniger genau
\end{tabular}
\end{table}

\begin{exampleblock}{Trade-off}
Genauigkeit vs. Performance \\
Für langsam bewegende Systeme: Höhere Intervalle OK
\end{exampleblock}
\end{frame}

%%%%%%%%%%%%%%%%%%%%%%%%%%%%%%%%%%%%%%%%%%%%%%%%%
\section{Demo \& Ergebnisse}
%%%%%%%%%%%%%%%%%%%%%%%%%%%%%%%%%%%%%%%%%%%%%%%%%

\begin{frame}{Live Simulationen - Demo}
\begin{center}
\vspace{2cm}
{\Huge \textbf{LIVE DEMO}}

\vspace{1cm}
{\Large Verschiedene Astrophysik-Szenarien}

\vspace{0.5cm}
\begin{itemize}
    \item Keplerian Disk
    \item Big Bang
    \item Globular Cluster
    \item Galaxy Collision
    \item Black Hole Accretion
    \item Brute Force vs. Barnes-Hut Vergleich
    \item Theta-Parameter live anpassen
\end{itemize}
\end{center}
\end{frame}

\begin{frame}{Performance-Vergleich}
\begin{block}{Eigene Messungen (1.000 Partikel)}
\begin{table}
\centering
\begin{tabular}{l|r|r}
\textbf{Methode} & \textbf{Berechnungen/Frame} & \textbf{Speedup} \\ \hline
Brute Force & 499.500 & 1x \\
Barnes-Hut ($\theta=0.5$) & $\sim$25.000 & 20x \\
Barnes-Hut ($\theta=0.7$) & $\sim$15.000 & \textbf{33x} \\
Barnes-Hut ($\theta=1.0$) & $\sim$10.000 & 50x
\end{tabular}
\end{table}
\end{block}

\begin{alertblock}{Ergebnis}
Barnes-Hut ist \textbf{30-50x schneller} bei guter Genauigkeit!
\end{alertblock}
\end{frame}

%%%%%%%%%%%%%%%%%%%%%%%%%%%%%%%%%%%%%%%%%%%%%%%%%
\section{Zusammenfassung}
%%%%%%%%%%%%%%%%%%%%%%%%%%%%%%%%%%%%%%%%%%%%%%%%%

\begin{frame}{Zusammenfassung}
\begin{block}{Barnes-Hut Algorithmus}
Hierarchische Approximationsmethode für N-Körper-Problem
\end{block}

\textbf{Kernkonzepte:}
\begin{enumerate}
    \item \textbf{Octree}: Hierarchische Raumpartitionierung
    \item \textbf{Massenschwerpunkt}: Bottom-up Berechnung
    \item \textbf{Theta-Kriterium}: Approximation vs. Genauigkeit
    \item \textbf{Komplexität}: $O(n^2) \rightarrow O(n \log n)$
\end{enumerate}

\vspace{0.3cm}
\textbf{Vorteile:}
\begin{itemize}
    \item 30-50x Speedup gegenüber Brute Force
    \item Skaliert gut mit N
    \item Anpassbare Genauigkeit ($\theta$-Parameter)
\end{itemize}
\end{frame}

\begin{frame}{Quellen}
\begin{thebibliography}{99}
\bibitem{barnes1986}
Josh Barnes, Piet Hut (1986):
\textit{A Hierarchical O(N log N) Force-Calculation Algorithm},
Nature, Vol. 324, pp. 446-449

\bibitem{wiki-octree}
Wikipedia: Octree Data Structure \\
\url{https://en.wikipedia.org/wiki/Octree}

\bibitem{wiki-nbody}
Wikipedia: N-body Problem \\
\url{https://en.wikipedia.org/wiki/N-body_problem}

\bibitem{raylib}
Raylib Documentation (Visualisierung) \\
\url{https://www.raylib.com/}
\end{thebibliography}

\vspace{0.3cm}
\begin{center}
\textbf{Eigene C++ Implementierung auf GitHub} \\
\textit{[Link einfügen]}
\end{center}
\end{frame}

\begin{frame}
\begin{center}
{\Huge Vielen Dank für Ihre Aufmerksamkeit!}

\vspace{1cm}
{\Large Fragen?}
\end{center}
\end{frame}

\end{document}
